% GENERATED FILE. Edit canonical lessons, then run npm run generate:books.
% canonical-source-hash: fnv1a64:9c0d473bb17a1396
% canonical-lessons: TA-C306-talaippuc-ceyti, TA-C306-mutal-pakkam, TA-C306-ital, TA-C306-patippu, TA-C306-veliyitu

\chapter{Front Page News}
\label{ch:ta-front-page-news}
\hlchaptermodality{306}

\begin{chapteropening}
\textbf{By the end of this chapter:} \emph{I can talk about newspapers and what they print.}

Complete the last lesson of chapter 306: I can talk about newspapers and what they print.
\end{chapteropening}

\section[\texorpdfstring{\ta{தலைப்புச்} \ta{செய்தி} (talaippuc ceyti)}{தலைப்புச் செய்தி (talaippuc ceyti)}]{\ta{தலைப்புச்} \ta{செய்தி} (talaippuc ceyti) — a headline}

\label{lesson:TA-C306-talaippuc-ceyti}



\begin{quote}
Before the new one: say the Tamil for a presenter, then the Tamil for a cameraman.
\end{quote}

\subsection*{You'll want to know: \ta{தலைப்புச்} \ta{செய்தி}}
\textbf{\ta{தலைப்புச்} \ta{செய்தி}} — \emph{talaippuc ceyti} — \textquotedblleft{}a headline\textquotedblright{}.

\begin{quote}
\ta{இன்றைய} \ta{தலைப்புச்} \ta{செய்தி} \ta{அந்த} \ta{விபத்துதான்}.
\end{quote}

\begin{quote}
\emph{iṉṟaiya talaippuc ceyti anta vipattutāṉ.}
\end{quote}

\begin{quote}
Today's headline is that very accident.
\end{quote}

The \ta{ச்} is the doubled stop of sandhi before \ta{செய்தி}, news.

\begin{cousinweb}
\textbf{\ta{தலைப்புச்} \ta{செய்தி}} (\emph{talaippuc ceyti}) is a Tamil phrase built from \ta{தலைப்பு}, \textquotedblleft{}beginning, title\textquotedblright{}, from \ta{தலை}, \textquotedblleft{}head\textquotedblright{} (DEDR 3103), and \ta{செய்தி}, \textquotedblleft{}news\textquotedblright{}, at root \textquotedblleft{}a deed\textquotedblright{} (DEDR 1957).
\end{cousinweb}

\begin{grammarlens}[title={What you've built}]
One more word for telling what happened.
\end{grammarlens}

\subsection*{Your turn}
\begin{itemize}
\raggedright
  \item \emph{Say it:} \emph{talaippuc ceyti}
  \item \emph{Say it:} \emph{talaippuc ceyti}, once more
  \item \emph{Say it:} \emph{iṉṟaiya talaippuc ceyti anta vipattutāṉ.} — Today's headline is that very accident.
  \item \emph{Recall:} say the Tamil for a presenter, then the Tamil for a cameraman, then say \emph{talaippuc ceyti} again
\end{itemize}

\begin{tcolorbox}[breakable,colback=teal!4,colframe=teal!35!black,title={Before you move on}]
What does \ta{தலைப்புச்} \ta{செய்தி} mean? (\textquotedblleft{}A headline\textquotedblright{}.) Say it once more.
\end{tcolorbox}
\section[\texorpdfstring{\ta{முதல்} \ta{பக்கம்} (mutal pakkam)}{முதல் பக்கம் (mutal pakkam)}]{\ta{முதல்} \ta{பக்கம்} (mutal pakkam) — the front page}

\label{lesson:TA-C306-mutal-pakkam}



\begin{quote}
Before the new one: say the Tamil for a cameraman, then the Tamil for a headline.
\end{quote}

\subsection*{You'll want to know: \ta{முதல்} \ta{பக்கம்}}
\textbf{\ta{முதல்} \ta{பக்கம்}} — \emph{mutal pakkam} — \textquotedblleft{}the front page\textquotedblright{}.

\begin{quote}
\ta{அவருடைய} \ta{படம்} \ta{முதல்} \ta{பக்கத்தில்} \ta{வந்தது}.
\end{quote}

\begin{quote}
\emph{avaruṭaiya paṭam mutal pakkattil vantatu.}
\end{quote}

\begin{quote}
His photo appeared on the front page.
\end{quote}

\ta{முதல்} \ta{பக்கத்தில்}, on the front page.

\begin{cousinweb}
\textbf{\ta{முதல்} \ta{பக்கம்}} (\emph{mutal pakkam}) joins native \ta{முதல்}, \textquotedblleft{}first\textquotedblright{} (DEDR 4950), and \ta{பக்கம்}, \textquotedblleft{}side, page\textquotedblright{}, borrowed from Sanskrit \emph{pakṣa}, \textquotedblleft{}wing, side\textquotedblright{} (CDIAL 7627).
\end{cousinweb}

\begin{grammarlens}[title={What you've built}]
One more word for telling what happened.
\end{grammarlens}

\subsection*{Your turn}
\begin{itemize}
\raggedright
  \item \emph{Say it:} \emph{mutal pakkam}
  \item \emph{Say it:} \emph{mutal pakkam}, once more
  \item \emph{Say it:} \emph{avaruṭaiya paṭam mutal pakkattil vantatu.} — His photo appeared on the front page.
  \item \emph{Recall:} say the Tamil for a cameraman, then the Tamil for a headline, then say \emph{mutal pakkam} again
\end{itemize}

\begin{tcolorbox}[breakable,colback=teal!4,colframe=teal!35!black,title={Before you move on}]
What does \ta{முதல்} \ta{பக்கம்} mean? (\textquotedblleft{}The front page\textquotedblright{}.) Say it once more.
\end{tcolorbox}
\section[\texorpdfstring{\ta{இதழ்} (itaḻ)}{இதழ் (itaḻ)}]{\ta{இதழ்} (itaḻ) — an issue; a magazine}

\label{lesson:TA-C306-ital}



\begin{quote}
Before the new one: say the Tamil for a headline, then the Tamil for the front page.
\end{quote}

\subsection*{You'll want to know: \ta{இதழ்}}
\textbf{\ta{இதழ்}} — \emph{itaḻ} — \textquotedblleft{}an issue; a magazine\textquotedblright{}.

\begin{quote}
\ta{இந்த} \ta{வார} \ta{இதழில்} \ta{ஒரு} \ta{நல்ல} \ta{கதை} \ta{வந்தது}.
\end{quote}

\begin{quote}
\emph{inta vāra itaḻil oru nalla katai vantatu.}
\end{quote}

\begin{quote}
A good story appeared in this week's issue.
\end{quote}

In older Tamil \ta{இதழ்} is a petal or a palm leaf; \ta{நாளிதழ்}, a daily paper, is \ta{நாள்} + \ta{இதழ்}.

\begin{cousinweb}
\textbf{\ta{இதழ்}} (\emph{itaḻ}) is native Dravidian (DEDR 459): a petal or a palmyra leaf, the leaf people once wrote on, and so an issue of a paper.
\end{cousinweb}

\begin{grammarlens}[title={What you've built}]
One more word for telling what happened.
\end{grammarlens}

\subsection*{Your turn}
\begin{itemize}
\raggedright
  \item \emph{Say it:} \emph{itaḻ}
  \item \emph{Say it:} \emph{itaḻ}, once more
  \item \emph{Say it:} \emph{inta vāra itaḻil oru nalla katai vantatu.} — A good story appeared in this week's issue.
  \item \emph{Recall:} say the Tamil for a headline, then the Tamil for the front page, then say \emph{itaḻ} again
\end{itemize}

\begin{tcolorbox}[breakable,colback=teal!4,colframe=teal!35!black,title={Before you move on}]
What does \ta{இதழ்} mean? (\textquotedblleft{}An issue; a magazine\textquotedblright{}.) Say it once more.
\end{tcolorbox}
\section[\texorpdfstring{\ta{பதிப்பு} (patippu)}{பதிப்பு (patippu)}]{\ta{பதிப்பு} (patippu) — an edition}

\label{lesson:TA-C306-patippu}



\begin{quote}
Before the new one: say the Tamil for the front page, then the Tamil for an issue.
\end{quote}

\subsection*{You'll want to know: \ta{பதிப்பு}}
\textbf{\ta{பதிப்பு}} — \emph{patippu} — \textquotedblleft{}an edition\textquotedblright{}.

\begin{quote}
\ta{புத்தகத்தின்} \ta{இரண்டாம்} \ta{பதிப்பு} \ta{வெளிவந்தது}.
\end{quote}

\begin{quote}
\emph{puttakattiṉ iraṇṭām patippu veḷivantatu.}
\end{quote}

\begin{quote}
The book's second edition came out.
\end{quote}

From \ta{பதி}, to imprint.

\begin{cousinweb}
\textbf{\ta{பதிப்பு}} (\emph{patippu}) is native Dravidian (DEDR 3911): from \ta{பதி}, \textquotedblleft{}to be imprinted\textquotedblright{}, an imprinting, and so a printing.
\end{cousinweb}

\begin{grammarlens}[title={What you've built}]
One more word for telling what happened.
\end{grammarlens}

\subsection*{Your turn}
\begin{itemize}
\raggedright
  \item \emph{Say it:} \emph{patippu}
  \item \emph{Say it:} \emph{patippu}, once more
  \item \emph{Say it:} \emph{puttakattiṉ iraṇṭām patippu veḷivantatu.} — The book's second edition came out.
  \item \emph{Recall:} say the Tamil for the front page, then the Tamil for an issue, then say \emph{patippu} again
\end{itemize}

\begin{tcolorbox}[breakable,colback=teal!4,colframe=teal!35!black,title={Before you move on}]
What does \ta{பதிப்பு} mean? (\textquotedblleft{}An edition\textquotedblright{}.) Say it once more.
\end{tcolorbox}
\section[\texorpdfstring{\ta{வெளியிடு} (veḷiyiṭu)}{வெளியிடு (veḷiyiṭu)}]{\ta{வெளியிடு} (veḷiyiṭu) — to bring out; to publish; to release}

\label{lesson:TA-C306-veliyitu}



\begin{quote}
Before the new one: say the Tamil for an issue, then the Tamil for an edition.
\end{quote}

\subsection*{You'll want to know: \ta{வெளியிடு}}
\textbf{\ta{வெளியிடு}} — \emph{veḷiyiṭu} — \textquotedblleft{}to bring out; to publish; to release\textquotedblright{}.

\begin{quote}
\ta{நாளிதழ்} \ta{அவருடைய} \ta{கடிதத்தை} \ta{வெளியிட்டது}.
\end{quote}

\begin{quote}
\emph{nāḷitaḻ avaruṭaiya kaṭitattai veḷiyiṭṭatu.}
\end{quote}

\begin{quote}
The newspaper published his letter.
\end{quote}

Past (I …): \ta{வெளியிட்டேன்}, \emph{veḷiyiṭṭēṉ}.

\ta{வெளி}, outside + \ta{இடு}, to put: to put out.

\begin{cousinweb}
\textbf{\ta{வெளியிடு}} (\emph{veḷiyiṭu}) is a Tamil compound built from \ta{வெளி}, \textquotedblleft{}outside, the open, public\textquotedblright{} (DEDR 5498), and \ta{இடு}, \textquotedblleft{}to put\textquotedblright{}.
\end{cousinweb}

\begin{grammarlens}[title={What you've built}]
One more word for telling what happened.
\end{grammarlens}

\subsection*{Your turn}
\begin{itemize}
\raggedright
  \item \emph{Say it:} \emph{veḷiyiṭu}
  \item \emph{Say it:} \emph{veḷiyiṭu}, once more
  \item \emph{Say it:} \emph{nāḷitaḻ avaruṭaiya kaṭitattai veḷiyiṭṭatu.} — The newspaper published his letter.
  \item \emph{Recall:} say the Tamil for an issue, then the Tamil for an edition, then say \emph{veḷiyiṭu} again
\end{itemize}

\begin{tcolorbox}[breakable,colback=teal!4,colframe=teal!35!black,title={Before you move on}]
What does \ta{வெளியிடு} mean? (\textquotedblleft{}To bring out; to publish; to release\textquotedblright{}.) Say it once more.
\end{tcolorbox}
